\subsection{BÀI TẬP TRẮC NGHIỆM}
\setcounter{ex}{0}
%\columnsep=25pt
%\begin{multicols}{2}
\ind{PHẦN I.} \inden{Câu trắc nghiệm nhiều phương án lựa chọn. Học sinh trả lời từ câu 1 đến câu 16. Mỗi câu hỏi học sinh chỉ chọn một phương án.}\\
\setcounter{ex}{0}
\Opensolutionfile{ans}[ans/1D2-B2-1]
\begin{ex}%[Học kỳ 1 lớp 11, THPT Lý Thái Tổ - Bắc Ninh, 2018-2019]%[Đình Nguyên, 11EX-HK1-1819]%[1D3Y3-1]
	Trong các dãy số hữu hạn sau, dãy số nào là cấp số cộng?
	\choice
	{$2$; $8$; $32$}
	{$3$; $7$; $11$; $16$}
	{\True $(u_n)$ với $u_n=4+3n$}
	{$(v_n)$ với $v_n=n^3$}
	\loigiai{
		Ta có $u_{n+1}-u_n=(4+3(n+1))-(4+3n)=3$. Suy ra dãy số $(u_n)$ là cấp số cộng với công sai $d=3$.	
	}
\end{ex}

\begin{ex}
	Cho cấp số cộng $\left(u_n\right)$ biết $u_n=3-5n$. Tìm công sai $d$ của cấp số cộng $\left(u_n\right)$.
	\choice
	{$d=3$}
	{\True $d=-5$}
	{$d=-3$}
	{$d=5$}
	\loigiai{
		
		Theo giả thiết suy ra $ u_{n+1}=3-5(n+1)=-2-5n\Rightarrow u_{n+1}-u_n=-5, \forall n\ge 1.$
		
		Suy ra $ (u_n) $ là cấp số cộng, công sai $ d=-5. $
	}
\end{ex}

\begin{ex}
	Tìm giá trị của $x$, $y$ sao cho dãy số $-2, x, 4, y$ theo thứ tự lập thành một cấp số cộng.
	\choice
	{$x=2,y=8$}
	{\True $x=1,y=7$}
	{$x=2,y=10$}
	{$x=-6,y=2$}
	\loigiai{
		Dãy số $-2, x, 4, y$ theo thứ tự lập thành một cấp số cộng $\Rightarrow\left\{\begin{aligned}x&=\dfrac{-2+4}{2}\\4&=\dfrac{x+y}{2}\end{aligned}\right.\Leftrightarrow\left\{\begin{aligned}x&=1\\y&=7\end{aligned}\right.$.
	}
\end{ex}

\begin{ex}%[Thi học kỳ 1, Dĩ An, Bình Dương, 2018]%[1D3B3]
	Cho $9$, $x$, $-1$, $y$ là $4$ số lập thành cấp số cộng, khi đó giá trị của $x,y$ là  
	\choice
	{\True $\heva{&x=4\\&y=-6}$}
	{$\heva{&x=2\\&y=-6}$}
	{$\heva{&x=2\\&y=5}$}
	{$\heva{&x=4\\&y=6}$}
	\loigiai{
		Vì các số $9$, $x$, $-1$, $y$ lập thành cấp số cộng nên ta có :\\
		$x=\dfrac{9+(-1)}{2}=4$, $\dfrac{x+y}{2}=-1\Rightarrow y=-2-x=-6$.
	}
\end{ex}

\begin{ex}
	Tìm số hạng thứ $11$ của cấp số cộng có số hạng đầu bằng $3$ và công sai $d=-2$.
	\choice
	{$-21$}
	{$23$}
	{\True $-17$}
	{$-19$}
	\loigiai{
		$u_{11} = u_1 + 10d = -17$.
	}
\end{ex}

\begin{ex}%[Học kì 1 Toán 11, 2017 - 2018, Trấn Biên, Đồng Nai]%[1D3B3]
	Cho hai số $-3$ và $23$. Xen kẽ giữa hai số đã cho $n$ số hạng để tất cả các số đó tạo thành cấp số cộng có công sai $d=2$. Tìm $n$.
	\choice
	{$n = 14$}
	{$n = 15$}
	{$n = 13$}
	{\True $n = 12$}
	\loigiai{
		Cấp số cộng có $k$ số hạng gồm có $u_1=-3$ và số hạng cuối $u_k=23$.\\
		Ta có $u_k=u_1+\left(k-1\right)d \Leftrightarrow 23=-3+2\left(k-1\right) \Leftrightarrow k=14$.
		Do đó $n=k-2=12$.
	}
\end{ex}

\begin{ex}%[Đề KSCL HK1 Toán 11, B Bình Lục, Hà Nam, 2018]%[1D3B3]
	Cấp số cộng $(u_n)$ có $u_6=12$, $u_{10}=24$. Tìm số hạng đầu $u_1$.
	\choice
	{$u_1=3$}
	{$u_1=2$}
	{$u_1=5$}
	{\True $u_1=-3$}
	\loigiai{Giả sử cấp số cộng có công sai là $d$. Ta có hệ
		\begin{center}
			$\heva{&u_6=u_1+5d\\&u_{10}=u_1+9d}\Leftrightarrow\heva{&12=u_1+5d\\&24=u_1+9d}\Leftrightarrow\heva{&u_1=-3\\&d=3}.$
		\end{center}
	}
\end{ex}

\begin{ex}
	Cho cấp số cộng $ (u_{n}) $ với $ u_{1}=2 $, $ d=9 $. Khi đó số $ 2018 $ là số hạng thứ mấy trong dãy?
	\choice
	{$ 223 $}
	{\True $ 225 $}
	{$ 224 $}
	{$ 226 $}
	\loigiai{
		Ta có $ u_{n}=u_{1}+(n-1)d\Leftrightarrow 2018=2+(n-1)\cdot 9\Leftrightarrow n=225 $.
	}
\end{ex}

\begin{ex} %[HK1, Phan Bội Châu Đắk Lắk, 2018]%[1D3B3]
	Cho dãy số $\left({u_n}\right)$, biết: $u_1=3,u_{n+1}=u_n+4$ với $n\geqslant 1$. Tìm $u_{1000}$.
	\choice
	{$3900$}
	{$4000$}
	{\True $3999$}
	{$4200$}
	\loigiai{
		Dãy $\left(u_n\right)$ là dãy CSC có $u_1=3$ và $d=4$.\\
		$u_{n}==u_1+\left(n-1\right)d=3+4\left(n-1\right)$.\\
		Vậy $u_{1000}=3+4\cdot999=3999$.
	}
\end{ex}

\begin{ex}%[Thi HK1, THPT Lương Thế Vinh Hà Nội, 2019]%[Tuan Nguyen, dự án EX11]%[1D3B3-3]
	Cho cấp số cộng $(u_n)$ thỏa mãn $\heva{&u_4=7u_1\\&S_{5}=75}$. Tìm số hạng thứ hai của cấp số cộng này.
	\choice
	{\True $u_2=9$}
	{$u_2=6$}
	{$ u_2=3$}
	{$ u_2=12$}
	\loigiai{
		Từ $u_4=7u_1$ suy ra $u_1+3d=7u_1\Leftrightarrow d=2u_1$.\\
		Mà $S_5=75\Leftrightarrow \dfrac{5(2u_1+4d)}{2}=75\Leftrightarrow 25u_1=75\Leftrightarrow u_1=3, d=6$.\\
		Do đó $u_2=9$.
	}
\end{ex}



\begin{ex}%[1D3K3]
	Tính số hạng đầu $u_1$ và và công sai $d$ của cấp số cộng $\left(u_n\right)$, biết $\heva{
		& u_1+u_5-u_3=10 \\ 
		& u_1+u_6=7. \\}$
	\choice
	{$u_1=-36, d=13$}
	{$u_1=36, d=13$}
	{\True $u_1=36, d=-13$}
	{$u_1=-36, d=-13$}
	\loigiai{
		
		Theo giả thiết suy ra $ \heva{&u_1+2d=10\\&2u_1+5d=7}\Leftrightarrow \heva{&u_1=36\\&d=-13.} $
	}
\end{ex}

\begin{ex}%[1D3B3-5]%[Phan Tấn Phú, 11-HK2]
	Một cấp số cộng có $15$ số hạng. Biết tổng của $15$ số hạng đó bằng $120$ và công sai bằng $-4$. Tìm số hạng đầu.
	\choice
	{$u_1=-20$}
	{\True $u_1=36$}
	{$u_1=540$}
	{$u_1=64$}
	\loigiai{
		Từ công thức $S_n=\dfrac{n}{2}\left[2u_1+(n-1)d\right]$.\\
		thay $n=15$, $S_{15}=120$, $d=-4$ ta có $120=\dfrac{15}{2}\left[2u_1+14.(-4)\right]$. Từ đó tìm được $u_1=36$.
		
	}
\end{ex}

\begin{ex}%[HK1, Chuyên Lê Khiết, Q-Ngãi]%[1D3B3]
	Cho cấp số cộng $(u_n)$ có $u_1=3$ và công sai $d=-2$. Tính $S_{2017}=u_1+u_2+\cdots +u_{2017}$.
	\choice
	{$S_{2017}=-4060211$}
	{\True $S_{2017}=-4060221$}
	{$S_{2017}=4072323$}
	{$S_{2017}=4073232$}
	\loigiai{
		$S_n=u_1\cdot n+\dfrac{n(n-1)}{2}\cdot d\Rightarrow S_{2017}=3\cdot 2017+\dfrac{2017\cdot 2016}{2}\cdot (-2)=-4060221$.
	}
\end{ex}

\begin{ex}%[1D3B3]
	Tính tổng $S=1+5+9+\cdots+397$ ta được kết quả
	\choice
	{$19298$}
	{$19090$}
	{$19920$}
	{\True $19900$}
	\loigiai{
		$S=1+5+9+\cdots+397$ là tổng $100$ số hạng đầu tiên của cấp số cộng có $u_1=1$ và công sai $d=4$.\\
		$S_{100}=\dfrac{100}{2}\left(2\cdot1+99\cdot4\right)=19900$.
	}
\end{ex}

\begin{ex}%[1D3B3-6]
	Ông X vay của công ty A một khoản tiền $ 72 $ triệu đồng và ông này trả nợ cho công ty A như sau:
	Ở quý thứ nhất ông trả $ 3 $ triệu đồng và kể từ quý thứ $ 2 $ mức trả sẽ tăng thêm $ 0{,}2 $ triệu đồng mỗi quý.
	Hỏi sau bao lâu thì ông X trả hết nợ?
	\choice
	{$ 5 $ năm}
	{$ 6 $ năm}
	{$ 3 $ năm}
	{\True $ 4 $ năm}
	\loigiai{
		Gọi $ n $ là số quý mà ông X trả hết nợ.\\
		Gọi số tiền ông X trả mỗi quý lần lượt là $ u_1 $, $ u_2 $, $\ldots$, $ u_n $. Và $ (u_n) $ tạo thành cấp số cộng với $ u_1=3 $, công sai $ d=0{,}2 $. \\
		Khi đó $ u_1+u_2+\ldots +u_n =72 \Leftrightarrow n\cdot 3+\dfrac{n(n-1) \cdot 0{,}2}{2}=72 \Leftrightarrow \hoac{&n=16\\&n=-45 \text{ (loại)}.} $\\
		Tức là trong khoảng $ 4 $ (năm) thì ông X trả hết nợ.
	}
\end{ex}

\begin{ex}%[1D3B3-6]%
	Một đa giác có $n$ cạnh và có chu vi bằng $158$ cm. Biết số đo các cạnh của đa giác lập thành một cấp số cộng và công sai $d=3$ cm và cạnh lớn nhất có độ dài là $44$ cm. Đa giác có số cạnh $n$ bằng
	\choice
	{$n=5$}
	{$n=7$}
	{$n=6$}
	{\True $n=4$}
	\loigiai{
		Gọi $u_i$ ($1\leqslant i \leqslant n$) là độ dài cạnh thứ $i$ trong cấp số cộng tạo bởi đồ dài các cạnh của đa giác.\\
		Ta có $u_n=44$, $d=3$ nên $u_1=u_n-(n-1)d=44-(n-1)\cdot 3=47-3n$.\\
		Chu vi của đa giác là tổng các cạnh của đa giác và bằng $$\dfrac{n(u_1+u_n)}{2}=158\Leftrightarrow n(91-3n)=316\Leftrightarrow 3n^2-91n+316=0\Leftrightarrow \hoac{& n=4\\ & n=\dfrac{79}{3}\ (\text{loại}).}$$
		Vậy $n=4$ thỏa mãn yêu cầu bài toán.
	}
\end{ex}

\Closesolutionfile{ans}

\ind{PHẦN II.} \inden{Câu trắc nghiệm đúng sai. Học sinh trả lời từ câu 17 đến câu 20. Trong mỗi ý a), b), c), d) ở mỗi câu, học sinh chọn đúng hoặc sai.}\\
\Opensolutionfile{ans}[ans/1D2-B2-2]

\begin{ex}%[1D2B2-1]%[1D2B2-3]%[1D2B2-5]
	Cho cấp số cộng $(u_n)$ có số hạng tổng quát $u_n=7n-3$. Xét tính đúng sai của các khẳng định sau:
	\choiceTF
	{\True Số hạng đầu của cấp số cộng $(u_n)$ là $u_1=4$}
	{Công sai của cấp số cộng $(u_n)$ là $d=-3$}
	{\True $u_{2012}=14081$}
	{$S_{100}=697$}
	\loigiai{
		\begin{itemchoice}
			\itemch Từ công thức số hạng tổng quát $u_n=7n-3$ ta có $u_1=7 \cdot 1-3=4$.
			\itemch Với mọi $n \in \mathbb{N}^*$ ta có $d=u_{n+1}-u_n=7(n+1)-3-(7n-3)=7$.
			\itemch Ta có $u_{2012}=7 \cdot 2012-3=14081$.
			\itemch Ta có cấp số cộng $(u_n)$ có số hạng đầu $u_1=4$ và công sai $d=7$. Suy ra 
			$$S_n=\dfrac{n.[2u_1+(n-1) \cdot d]}{2}$$ nên $$S_{100}=\dfrac{100 \cdot [2\cdot 4+(100-1)\cdot 7]}{2}=35050.$$
		\end{itemchoice}
	}
\end{ex}

\begin{ex}%[1D2B2-2]
	Cho cấp số cộng $(u_n)$ biết $u_1=5$ và $d=3$. Xét tính đúng sai của các khẳng định sau:
	\choiceTF
	{\True $u_{99}=299$}
	{Số hạng tổng quát của cấp số cộng $(u_n)$ là $u_n=3n+5$}
	{\True Số $1502$ là số hạng thứ $500$ của cấp số cộng $(u_n)$}
	{\True Tổng của $150$ số hạng đầu bằng $34275$}
	\loigiai{
		\begin{itemchoice}
			\itemch Ta có $u_{99}=u_1+98\cdot d=5+98\cdot 3=299$.
			\itemch Số hạng tổng quát của cấp số cộng là \[u_n=u_1+(n-1).d=5+(n-1)\cdot 3=3n+2.\]
			\itemch Ta có \\
			$	\begin{aligned}
				u_n=u_1+(n-1).d &\Leftrightarrow 1502=5+3\cdot (n-1)\\
				&\Leftrightarrow 3n+2=1502\\
				&\Leftrightarrow n=500.
			\end{aligned}$\\	
			Vậy $u_{500}=1502$.
			\itemch Ta có $S_{150}=\dfrac{150}{2}\left(2u_1+149d \right)=34275.$
		\end{itemchoice}
		}
\end{ex}

\begin{ex}%[1C2K2-5]
	Cho $(u_n)$ là cấp số cộng có $u_2+u_4=22$, $u_1\cdot u_5=21$ và công sai $d$ dương. Xét tính đúng sai của các khẳng định sau:
	\choiceTF
	{\True Số hạng đầu là $u_1=1$}
	{Công sai là $d=4$}
	{\True $S_{100}=24850$}
	{$u_1+u_5+u_9+\cdots+u_{101}=6520$}
	\loigiai{
		Xét hệ $\heva{&u_2+u_4=22\\&u_1\cdot u_5=21} \Leftrightarrow \heva{&2u_1+4d=22\\&u_1\cdot \left(u_1+4d \right) =21} \Leftrightarrow \heva{&u_1=1\\&d=5}$
		\begin{itemchoice}
			\itemch Số hạng đầu $u_1=1$.
			\itemch Công sai $d=5$.
			\itemch $S_{100}=\dfrac{100}{2}\left(2u_1+99d \right)= 24850$.
			\itemch Các số $u_1$, $u_5$, $u_9$,$\ldots$, $u_{101}$ lập thành cấp số cộng có số hạng đầu $u_1=1$, công sai $d'=4d=20$.\\
			Tổng $u_1+u_5+u_9+\cdots+u_{101}$ gồm $26$ số hạng và bằng $6526$.
		\end{itemchoice}
	}
\end{ex}

\begin{ex}
	Khi kí kết hợp đồng lao động với người lao động, một doanh nghiệp đề xuất hai phương án trả lương như sau:
	\begin{itemize}
		\item Phương án 1: Năm thứ nhất, tiền lương là $ 120 $ triệu. Kể từ năm thứ hai trở đi, mỗi năm tiền lương được tăng $ 18 $ triệu.
		\item Phuơng án 2: Quý thứ nhất, tiền lương là $ 24 $ triệu. Kể từ quý thứ hai trở đi, mỗi quý tiền lương được tăng $ 1,8 $ triệu.
	\end{itemize}
	Xét tính đúng sai của các khẳng định sau:
	\choiceTF
	{Theo phương án 1, tổng số tiền người lao động nhận được sau 3 năm là 414 triệu đồng}
	{Theo phương án 2, tổng số tiền người lao động nhận được sau 3 năm là 416 triệu đồng}
	{Nếu kí hợp đồng lao động $ 3 $ năm, thì phương án 2 có lợi hơn}
	{Nếu kí hợp đồng lao động $ 10 $ năm, thì phương án 1 có lợi hơn}
	\loigiai{
	Ở phương án trả lương thứ nhất, số tiền lương mỗi năm người lao động nhận được lập thành cấp số cộng $(u_n)$ có số hạng đầu $u_1=120$, công sai $d=18$. \\
	Ở phương án trả lương thứ hai, số tiền lương mỗi quý người lao động nhận được lập thành cấp số cộng $\left(v_n\right)$ có số hạng đầu $v_1=24$, công sai $d'=1,8$.
	\begin{itemchoice}
		\itemch Tổng số tiền lương người lao động nhận được trong $3$ năm ở phương án $1$ là tổng $3$ số hạng đầu của cấp số cộng và bằng
		$$S_3=\dfrac{(2u_1+2d)\cdot 3}{2}=3u_1+3d=3\cdot 120+3\cdot 18=414 \text{ (triệu đồng).}$$
		\itemch Do $1$ năm có $4$ quý nên tổng số tiền lương người lao động nhận được trong $3$ năm ở phương án $2$ là tổng $12$ số hạng đầu của cấp số cộng và bằng
		$$S'_{12}=\dfrac{(2v_1+11d')\cdot 12}{2}=12v_1+66d'=12\cdot 24+66\cdot 1,8=406,8 \text{ (triệu đồng).}$$
		\itemch Vậy nếu kí hợp đồng lao động $3$ năm thì số tiền người lao động nhận được phương án $1$ nhiều hơn.
		\itemch Nếu kí hợp đồng lao động $10$ năm thì: \\
		Tổng số tiền lương người lao động nhận được trong $10$ năm ở phương án $1$ bằng
		$$S_{10}=\dfrac{\left(2 u_1+9 d\right) .10}{2}=10 u_1+45 d=10 \cdot 120+45.18=2010 \text{ (triệu đồng).}$$
		Tổng số tiền lương người lao động nhận được trong $10$ năm ở phương án $2$ bằng
		$$S'_{40}=\dfrac{(2v_1+39d')\cdot 40}{2}=40v_1+780d'=40\cdot 24+780 \cdot 1,8=2364 \text { (triệu đồng).}$$
		Vậy nếu kí hợp đồng lao động $10$ năm thì số tiền nhân được phương án $2$ nhiều hơn.
\end{itemchoice}}
\end{ex}

\Closesolutionfile{ans}

\ind{PHẦN III.} \inden{Câu trắc nghiệm trả lời ngắn. Học sinh trả lời từ câu 21 đến câu 25 vào ô kết quả.}\\
\Opensolutionfile{ans}[ans/1D2-B2-3]

\begin{ex}%[Thi học kỳ 1, Dĩ An, Bình Dương, 2018]%[1D3B3]
	Cho cấp số cộng $\left(u_n\right)$ có $u_1=123$ và $u_3-u_{15}=84$. Số hạng $u_{17}$ bằng bao nhiêu?\\
	\shortans[oly]{$11$}	
	\loigiai{ Với $d$ là công bội của cấp số cộng ta có:\\
		$u_3-u_{15}=84\Leftrightarrow u_1+2d-u_1-14d=84\Leftrightarrow d=-7$. Từ đó suy ra, $u_{17}=u_1+16d=11$.
	}
\end{ex}

\begin{ex}%[1D3G3-6]
	Bác Bình muốn trồng cây cà phê trên một ngọn đồi như sau: Từ trên đỉnh đồi trồng hàng thứ nhất $2$ cây; đi xuống hàng thứ hai $5$ cây; đi xuống hàng thứ ba $8$ cây; $\ldots$; đi xuống hàng cuối cùng dưới chân đồi trồng $26$ cây. Theo cách trồng như trên bác bình trồng được bao nhiêu cây cà phê?\\
	\shortans[oly]{$126$}	
	\loigiai
	{
		Gọi $u_1;u_2;u_3; \ldots; u_n$ lần lượt là số cây hàng thứ nhất; hàng thứ hai; hàng thứ ba; \ldots; hàng thứ $n$.
		Ta có $u_1=2; u_2=5 \Rightarrow d=3$.\\
		Theo cách trồng của bác Bình số cây mỗi hàng lập thành một cấp số cộng với số hạng đầu $u_1=2$ và công sai $d=3$. Thật vậy, ta có:
		\begin{itemize}
			\item[•] $u_1= 2$ cây.
			\item[•] $u_2=u_1+d = 2+3=5$ cây.
			\item[•] $u_3=u_1+2d=5+3=8$ cây.
			\item[•] \ldots
			\item[•] $u_n = u_1+(n-1)d = 26$ cây.
		\end{itemize}
		Suy ra $2+(n-1)\cdot 3 =26 \Leftrightarrow n = 9$. Vậy có $9$ hàng cây.\\
		Do đó, số cây cà phê bác Bình trồng được là
		\[S_9 =\dfrac{9(2+26)}{2}=126 \text{ cây.} \]
	}
\end{ex}

\begin{ex}%[1D3B3-6]
	Một công ty trả lương cho anh A mức lương là $ 4{,}5 $ triệu đồng/quý và kể từ quý làm việc thứ 2 thì mức lương sẽ tăng thêm $ 0{,}3 $ triệu đồng mỗi quý. Hỏi tổng số tiền  (tính bằng triệu)  sau $ 3 $ năm làm việc anh A nhận được là bao nhiêu? \\
	\shortans[oly]{$73{,}8$}	
	\loigiai{
		$ 3 $ năm có $ 12 $ quý.\\
		Gọi số tiền Anh A lãnh ở quý $ n $ là $ u_n $. Khi đó $ (u_n) $ lập thành cấp số cộng có $ u_1=4{,}5 $ và công sai $ d=0{,3}  $.\\
		Tổng số tiền sau $ 3 $ năm anh A nhận được là
		\[ S_{12} =u_1+u_{2}+\cdots+u_{12}=12 \cdot 4{,}5+\dfrac{12 \cdot 11  \cdot 0{,}3}{2} = 73{,}8 \text{ (triệu).} \]
	}
\end{ex}


\begin{ex}%[1D3B3-6]%
	Bé An luyện tập khiêu vũ cho buổi dạ hội cuối khóa. Bé bắt đầu luyện tập trong $1$ giờ vào ngày đầu tiên. Mỗi ngày tiếp theo, bé tăng thêm $5$ phút luyện tập so với ngày trước đó. Hỏi sau một tuần, tổng thời gian bé An đã luyện tập là bao nhiêu phút?\\
	\shortans[oly]{$525$}	
	\loigiai{
		Thời gian luyện tập của bé An là một cấp số cộng với $u_1=60$ phút và $d=5$ phút.\\
		Vậy, sau một tuần tổng thời gian bé An luyện tập là $S_7=\dfrac{7}{2}\left(2\cdot 60+6\cdot 5\right)=525$ phút.
	}
\end{ex}

\begin{ex}
	Một cầu thang bằng gạch có tổng cộng $30$ bậc. Bậc dưới cùng cần $100$ viên gạch. Mỗi bậc tiếp theo cần ít hơn hai viên gạch so với bậc ngay trước đó. Hỏi cần bao nhiêu viên gạch để xây cầu thang?\\
	\shortans[oly]{$2130$}	
	\loigiai{
		Theo bài ra ta có số viên gạch ở mỗi bậc thang (tính từ dưới lên) lập thành một cấp số cộng gồm $30$ số với số hạng đầu $u_1=100$ và công sai $d=-22$.		\\
		Do đó, công thức của cấp số cộng biểu thị số viên gạch cho mỗi bậc cầu thang như sau:
		\[u_1=100;~u_{n+1}=u_n+(-2)\text{ với }n\ge 1.\]
		Ta có $S_{30}=u_1+u_2+\cdots +u_{30}=\dfrac{30}{2}\left[2\cdot 100+(30-1)(-2)\right]=2130$.\\
		Như vậy, ta cần $2 130$ viên gạch để xây cầu thang.
		}
\end{ex}

\Closesolutionfile{ans}

\Closesolutionfile{ans}

